\documentclass[11pt]{article}

\usepackage[T1]{fontenc}
\usepackage[margin=0.8in]{geometry}
\usepackage{enumitem}
\usepackage{xcolor}
\usepackage{hyperref}
\usepackage{parskip}
\usepackage{titlesec}

\definecolor{navy}{HTML}{102A43}
\definecolor{blue}{HTML}{1F5A94}
\definecolor{cyan}{HTML}{187E8A}
\definecolor{softgray}{HTML}{526575}

\hypersetup{
  colorlinks=true,
  linkcolor=blue,
  urlcolor=blue,
  pdftitle={Assured PNT, Part 1: Voice-Track Transcript},
  pdfauthor={Michael L. Carroll}
}
\titleformat{\section}{\Large\bfseries\color{navy}}{}{0pt}{}
\titleformat{\subsection}{\large\bfseries\color{blue}}{}{0pt}{}
\newcommand{\timing}[1]{\textcolor{softgray}{\textbf{Target: #1}}}
\newcommand{\cue}[1]{%
  \par\noindent\fcolorbox{cyan!70}{cyan!7}{%
    \parbox{0.94\linewidth}{\textbf{Production cue:} #1}}\par}
\newenvironment{narration}{\begin{quote}\itshape}{\end{quote}}

\title{Assured PNT, Part 1\\\large Navigation When GPS Cannot Be Trusted}
\author{Voice-track transcript for Michael L. Carroll}
\date{October 2026}

\begin{document}
\maketitle

\begin{center}
  \textbf{Recording target: 17--18 minutes} \quad $\vert$ \quad
  \textbf{15 logical slides / 44 build pages}
\end{center}

\section*{Recording notes}

This script contains about 1,800 spoken words. The 17-minute-20-second timing
plan assumes a deliberately measured technical delivery, including build cues
and pauses; a brisk continuous read will be closer to 14--15 minutes. Text in
the narration blocks is intended to be spoken. Production cues are not spoken.
A short pause means about one second; a full pause means two to three seconds.
Advance each build at the sentence indicated, leaving enough time for the
viewer to read the new element.

Pronunciations: PNT is ``pee-en-tee''; GNSS is ``gee-en-ess-ess''; INS is
``eye-en-ess''; IMU is ``eye-em-you''; NIS is ``en-eye-ess.'' Read $P$ as
``P,'' $R$ as ``R,'' and $\nu$ as ``nu.''

\section{Slide 1 --- Assured PNT}
\timing{0:25}

\cue{Open on the complete title slide. No build.}

\begin{narration}
What happens when a navigation system gives a precise, confident answer---and
that answer is wrong? In this first episode on assured position, navigation,
and timing, we will look at the system behavior required when GPS is blocked,
jammed, or deliberately manipulated. The central theme is simple: resilience
does not come from one perfect sensor. It comes from prediction, disciplined
trust, and deliberate recovery.
\end{narration}

\section{Slide 2 --- A confident answer can still be wrong}
\timing{1:15}

\cue{Begin on build 1. Advance once after ``those are different failures.''}

\begin{narration}
The most dangerous navigation failure is not always a blank screen. It is a
solution that still looks healthy after its reference has been manipulated.
The receiver may continue producing position, velocity, and time. The output
may be smooth. Its reported uncertainty may even remain small. Yet the answer
can be moving away from truth.

[Short pause.] A loss of availability says, ``I no longer have the
measurement.'' A loss of integrity says, ``I have a measurement, but I should
not rely on it.'' Those are different failures.

[Advance.] That distinction gives us three rules for the rest of the episode.
Availability is not integrity. Precision is not truth. And putting a
measurement into a fusion algorithm does not automatically make that
measurement trustworthy. A resilient system must keep its confidence tied to
evidence.
\end{narration}

\section{Slide 3 --- What PNT provides}
\timing{1:05}

\cue{Reveal Position, Navigation, and Timing one at a time. Reveal the summary last.}

\begin{narration}
[Build 1.] PNT begins with position: where am I?

[Build 2.] Navigation adds motion and intent: where am I going, how fast am I
moving, and how is the vehicle oriented?

[Build 3.] Timing asks when an observation or event occurred. Timing is not an
afterthought. It supports measurement alignment, communication, sensing, and
the clock states inside a GNSS solution.

[Build 4.] The important engineering question is not whether every one of
these quantities remains perfect under every condition. It is which
combination of position, velocity, attitude, and time the mission actually
needs---and what error it can tolerate. Assurance begins by turning that need
into explicit performance limits.
\end{narration}

\section{Slide 4 --- GPS is powerful---and physically fragile}
\timing{1:20}

\cue{Start with the left block. Reveal the right block, then the concluding line.}

\begin{narration}
GPS is an extraordinary navigation reference. It is global, common to many
users, and capable of providing position, velocity, and precise time with very
little infrastructure on the user platform. Its errors are also bounded over
the long term, which makes it an excellent partner for inertial navigation.

[Advance.] But the signal arriving at the antenna is weak. Buildings, terrain,
foliage, and vehicle structure can block it. Reflections can create multipath.
Radio-frequency interference can overpower it. And a counterfeit signal can
attempt to draw a receiver toward a false solution.

[Advance.] So our design objective cannot be nominal GPS accuracy alone. We
have to design for both loss and manipulation. That means preserving useful
navigation when measurements disappear, and protecting the estimator when
measurements remain present but become misleading.
\end{narration}

\section{Slide 5 --- Two failures, two different hazards}
\timing{1:25}

\cue{Show jamming first. Reveal spoofing second. Reveal the Kalman-filter caveat last.}

\begin{narration}
Jamming and spoofing create different estimator problems.

Jamming denies service. Interference raises the effective noise floor or
overwhelms the authentic satellite signal. A receiver may lose tracking, lose
satellites, or declare degraded performance. That is serious, but the failure
is often observable: the external updates stop, and the navigator must bridge
the gap.

[Advance.] Spoofing counterfeits the measurement. False GNSS-like data can
appear plausible and can pull the estimated state away from reality. If the
fusion system accepts that data repeatedly, it may also corrupt estimates of
inertial sensor bias or clock error. The damage can therefore persist after
the immediate event.

[Advance.] Statistical consistency checks can reject measurements that
disagree with prediction. But a Kalman filter, by itself, does not authenticate
a radio signal. A sophisticated false input can remain statistically plausible.
Integrity requires multiple layers of evidence.
\end{narration}

\section{Slide 6 --- Assurance is a behavior, not a sensor}
\timing{1:20}

\cue{Advance through Protect, Detect, Respond, and Recover; then reveal the descriptions and final pill.}

\begin{narration}
[Build 1.] Protect the system by reducing avoidable exposure: sound antenna
installation, interference awareness, secure interfaces, and sensible
operating procedures.

[Build 2.] Detect anomalies quickly. Monitor receiver health, signal
conditions, estimator consistency, and agreement among independent sources.

[Build 3.] Respond in a controlled way. Isolate a suspect input, change its
weight, switch modes, and continue only at a level of performance the system
can support.

[Build 4.] Recover deliberately. A source that has returned is not necessarily
a source that should be trusted immediately.

[Build 5.] These four actions connect technical design to operations: prevent
what we can, recognize what remains, degrade gracefully, and test sources
before re-admission.

[Build 6.] The system must know what it trusts---and why. That is behavior of
the complete architecture, not a property we can buy in a single sensor.
\end{narration}

\section{Slide 7 --- Inertial navigation provides the bridge}
\timing{1:20}

\cue{Begin with INS strengths. Reveal the uncertainty plot and limitation together.}

\begin{narration}
An inertial navigation system is the natural bridge through a GPS outage. Once
initialized, it is self-contained. Accelerometers and gyroscopes support a
high-rate estimate of position, velocity, and attitude. Because the measurement
is mechanical rather than radio-frequency based, an IMU is not disabled by an
RF jammer. And over short intervals, inertial propagation is smooth and
responsive.

[Advance.] The price is drift. Accelerometer bias is integrated into velocity
and then position. Gyroscope bias creates attitude error, which projects
gravity and specific force in the wrong direction. Noise and imperfect models
accumulate.

The curve is conceptual, not a universal error law. Its message is operational:
uncertainty grows with outage time. The vehicle may continue only until the
predicted uncertainty reaches the mission limit. INS gives us time, but it does
not give us infinite time.
\end{narration}

\section{Slide 8 --- Resilience comes from complementary errors}
\timing{1:15}

\cue{Show GNSS plus INS, add the other aids, then reveal the design test.}

\begin{narration}
[Build 1.] GPS and INS work well together because their errors are
complementary. Inertial sensors provide continuity and rapid dynamics but
drift. GNSS provides an external reference that can bound that drift, but its
signals can disappear or be deceived.

[Build 2.] The same reasoning extends to air data, barometric altitude,
terrestrial radio ranging, terrain matching, and vision. Each source observes
some part of the state, with its own geometry, update rate, environment, and
failure mechanisms.

[Build 3.] Simply adding sensors is not resilience. A useful source has failure
modes that are sufficiently independent of the sources already present. Its
measurements must also make the mission-critical states observable. Ten sensors
that share the same hidden dependency may provide less assurance than two
carefully chosen sources based on different physics.
\end{narration}

\section{Slide 9 --- The Kalman filter is the integration engine}
\timing{1:30}

\cue{The full diagram is visible. Trace it from left to right with the narration.}

\begin{narration}
The Kalman filter gives this architecture a disciplined predict-and-update
cycle.

IMU samples drive the prediction. The mechanization propagates the navigation
state, while the process model propagates the covariance, $P$. The superscript
minus means this is the prior estimate: our best prediction before using the
new external measurement.

External measurements then pass through validation. Does the observation agree
with what the predicted state and its uncertainty say is plausible? Do other
independent sources support it? Is the receiver healthy? If the evidence is
adequate, the update corrects the state and reduces uncertainty. The
superscript plus marks the posterior estimate.

If the evidence is poor, the system can reject the measurement, isolate the
source, or increase its assumed noise so that it receives less weight. The
important point is that prediction carries the solution through a gap, while
trusted updates bound drift. An update is beneficial only when accepting it is
safer than continuing without it.
\end{narration}

\section{Slide 10 --- A practical error-state model}
\timing{1:10}

\cue{Reveal the state vector, then the table, then the design-choice pill.}

\begin{narration}
[Build 1.] Many integrated navigators use an error-state filter. The inertial
mechanization carries the full navigation solution, and the Kalman filter
estimates small errors around that solution.

[Build 2.] Here the state contains position and velocity error, small attitude
error, accelerometer and gyro biases, and receiver clock bias and drift.
Estimating the inertial biases is especially important: an external reference
does more than correct the current position. It calibrates error sources that
would otherwise drive future drift.

[Build 3.] This vector is an example, not a universal recipe. The state is a
design choice. Add a quantity when the mission needs it, when a credible model
exists, and when the available measurements can actually observe it. A larger
state that cannot be supported by the measurement geometry can create false
confidence rather than capability.
\end{narration}

\section{Slide 11 --- Trust begins with the innovation}
\timing{1:30}

\cue{Reveal the innovation, then NIS, then the validation list, and finally the caveat.}

\begin{narration}
[Build 1.] The innovation is the measured value, $z$, minus the value predicted
from the prior state. In plain language: what did the sensor say that the
filter did not expect?

A large difference may indicate interference, a failed sensor, a bad model, or
an aggressive maneuver the estimator did not represent well. The residual is
evidence, not yet a diagnosis.

[Build 2.] The normalized innovation squared---the NIS---scales the residual by
its predicted covariance, $S$. This matters because ten meters may be alarming
in one geometry and entirely plausible in another. Under the filter's modeling
assumptions, the NIS can be compared with a statistical threshold.

[Build 3.] A practical validator does more. It checks persistence over time,
compares independent sources, considers signal and receiver health, and uses
explicit recovery logic.

[Build 4.] A gate tells us that a measurement is inconsistent with the current
model and evidence. It does not prove whether the cause was spoofing, jamming,
multipath, a sensor fault, or a bad assumption. Good assurance preserves that
distinction.
\end{narration}

\section{Slide 12 --- What the system should do during an outage}
\timing{1:25}

\cue{Advance through the timeline phases, then reveal the graceful-degradation block.}

\begin{narration}
[Build 1.] Before an event, trusted GNSS updates keep the inertial solution and
its bias estimates calibrated.

[Build 2.] When an anomaly appears, the system should detect it promptly and
stop treating the affected input as ordinary truth. That decision may be based
on several monitors rather than one threshold crossing.

[Build 3.] During the outage, the INS propagates the state. Other trusted aids
can continue updating the filter. Covariance, $P$, must grow consistently with
the remaining sensor errors and the weaker measurement geometry.

[Build 4.] When the source returns, validate it before reuse. Compare it against
the propagated solution and independent aids. Require a suitable interval of
consistent behavior. Re-admission may be gradual rather than instantaneous.

[Build 5.] This is graceful degradation: continue only while estimated
uncertainty remains inside a declared operational bound. When it does not,
transition to a safe mode. The estimate should never look healthier merely
because its best reference disappeared.
\end{narration}

\section{Slide 13 --- Four design rules for assured PNT}
\timing{1:15}

\cue{Reveal one numbered rule at a time, then reveal the final sentence.}

\begin{narration}
[Build 1.] First, diversify. Combine sources that use different physics and do
not share all of the same failure modes.

[Build 2.] Second, manage trust. A precise measurement is not necessarily an
authentic measurement. Weight must follow evidence.

[Build 3.] Third, expose uncertainty. Propagate covariance realistically,
validate it where possible, and compare it with the error the mission can
tolerate.

[Build 4.] Fourth, plan recovery before the failure occurs. Define how a source
is declared suspect, how it is isolated, what conditions permit continued
operation, and how it is tested before re-admission.

[Build 5.] These rules reinforce one another. Assurance is the controlled
behavior of the whole navigation system---including its estimator, monitors,
modes, interfaces, and operating procedures.
\end{narration}

\section{Slide 14 --- The central idea}
\timing{0:55}

\cue{Hold on the complete summary slide.}

\begin{narration}
The central idea can be reduced to three verbs: predict continuously,
challenge measurements, and recover deliberately.

A resilient navigator does not promise perfect availability. No physical
system can. It promises that performance is known, uncertainty is visible, and
behavior remains safe when important inputs are lost or challenged.

In Part Two, we will move from architecture to aiding sources. We will compare
barometric, radio, terrain, and vision-based measurements, and then examine how
terrain-referenced navigation can bound inertial drift when GNSS is not
available.
\end{narration}

\section{Slide 15 --- Primary references}
\timing{0:10}

\cue{Leave the references on screen for at least ten seconds. Do not narrate the individual URLs.}

\begin{narration}
The primary government references used for this episode are shown here and are
also linked in the video description.
\end{narration}

\section*{Source links for the video description}

\begin{itemize}[leftmargin=*]
  \item FAA, \emph{GPS and GNSS Interference Resource Guide}, 2025:\\
    \href{https://www.faa.gov/about/office_org/headquarters_offices/avs/offices/afx/afs400/afs410/GNSS}{FAA GPS/GNSS Interference Resource Guide page}
  \item GPS.gov, \emph{Resilience Through Responsible Use of PNT}:\\
    \url{https://www.gps.gov/resilience-through-responsible-use-pnt}
  \item NIST IR 8323 Rev. 1, \emph{Foundational PNT Profile}, 2023:\\
    \url{https://doi.org/10.6028/NIST.IR.8323r1}
  \item DHS, \emph{Resilient PNT Reference Architecture}, 2022:\\
    \url{https://www.dhs.gov/sites/default/files/2022-06/22_0609_st_resilient_pnt_ra.pdf}
  \item U.S. DOT, \emph{Complementary PNT Action Plan}, 2024:\\
    \url{https://www.transportation.gov/pnt/complementary-pnt-action-plan}
  \item DHS, \emph{Best Practices for Resilient PNT Supporting Critical
    Infrastructure}, 2025:\\
    \url{https://www.dhs.gov/sites/default/files/2025-02/25_0220_st_pnt_best_practices_ci.pdf}
\end{itemize}

\end{document}
